\scanpage{27}
\mspage{\handwritten{19}}
% The source is upside down; the reading was made from one 180-degree view
% of the whole and one 180-degree view of each supplied half.
% All text in the following figure belongs to the graphically cancelled draft.
% Local deletions retain their own deletion macros; the page-wide strokes are
% drawn across the complete text rather than replaced by horizontal strikeout.
\reversemarginpar
\figure{%
\begin{tikzpicture}
\node[anchor=north west,inner sep=0pt,outer sep=0pt] (draft) at (0,0) {%
\begin{minipage}{13cm}
\fontsize{9.5}{12.5}\selectfont
\setlength{\parindent}{0pt}
\csname tagsleft@true\endcsname
2.1.1. On peut dire encore qu'une structure de biextension sur le $G_{P\times Q}$-torseur $E$ est définie par deux lois de composition non partout définies $\handwritten{\phi}(x,y)$ et $\handwritten{\unsure{\mu|\psi}}(x,y)$ sur $E$, la première étant définie lorsque $\deleted{j(}$ $\mathrm{pr}_2j(x)=\mathrm{pr}_2j(y)$, la deuxième lorsque $\mathrm{pr}_1j(x)=\mathrm{pr}_1j(y)$ \deleted{(où} (où $j:E\handwritten{\longrightarrow}P\times Q$ est le morphisme structural du torseur). Les axiomes, en plus des associativités du type
\[
\deleted{\handwritten{\phi}(gx,y)\unsure{=}g\handwritten{\phi}(x,y),\quad\handwritten{\phi}(x,gy)\unsure{=}g\handwritten{\phi}(x,y),\quad\handwritten{\phi}(}
\]
\[
\handwritten{\phi}(gx,y)=\handwritten{\phi}(x,gy)=g\handwritten{\phi}(x,y)\qquad(x,y\handwritten{\in}E(S),\ g\handwritten{\in}G(S))
\tag*{(2.1.1.1)}
\]

et des relations analogues en $\handwritten{\mu}$, s'expriment alors par l\deleted{a}'associativité et la commutativité de \unsure{\deleted{s}c|\deleted{c}s}es lois de composition partielles, et par la relation
\[
\begin{tikzpicture}[baseline=(interchange.base),x=.008cm,y=-.008cm]
\node[anchor=base west,inner sep=0pt] (interchange) at (0,0) {$\displaystyle\handwritten{\phi}(\handwritten{\mu}(x,y),\handwritten{\mu}(z,t))=\handwritten{\mu}(\handwritten{\phi}(x,z),\handwritten{\phi}(y,t))$};
% The two distinct curves are kept with the expression they flank.
\draw[line width=.35pt] (-110,18) .. controls (-125,-7) and (-25,-23) .. (45,-31) .. controls (87,-41) and (133,-49) .. (151,-46) -- (154,-59);
\draw[line width=.4pt] (29,-52) .. controls (49,-32) and (-8,-17) .. (-20,-1) .. controls (-28,9) and (3,29) .. (19,33) -- (160,33) -- (162,47);
\end{tikzpicture}
\tag*{\raisebox{.33cm}{(2.1.1.2)}}
\]
pour $x,y,z,t$ respectivement dans $E_{p,q}(S),E_{p,q'}(S),E_{p',q}(S),E_{p',q'}(S)$.

\medskip
2.2. \inserted{\handwritten{\uline{Sections unité}.}} Il y a lieu de considérer la section unité de $\mathbb{E}\deleted{t}$ considéré comme extension de $Q_P$ par $G_P$,
\[
e_{E/P}\handwritten{\in\Gamma}(E/P)=\mathrm{Hom}_P(P,E)\ ,
\tag*{(2.2.1)}
\]
et de façon symétrique la section unité de $E$ considéré comme extension de $P_Q$ par $G_Q$,
\[
e_{E/Q}\handwritten{\in\Gamma}(E/Q)=\mathrm{Hom}_Q(Q,E)\ .
\tag*{(2.2.2)}
\]

En particulier, \uline{le torseur} $E$ \uline{au dessus de} $P\times Q$ (de \unsure{\deleted{g}G|\deleted{G}g}roupe $G_{P\times Q}$) \uline{est canoniquement trivialisé au dessus de} $1_P\times Q$ \uline{par} $e_{E/P}$, \uline{et au dessus de} $P\times1_Q$ \uline{par} $e_{E/Q}$.
\inserted{\begin{tikzpicture}[baseline=(deplus.base)]
\node[anchor=base west,inner sep=0pt] (deplus) at (0,.25) {De plus,};
\node[anchor=base west,inner sep=0pt] at (0,0) {\deleted{\handwritten{\illegible{1}}}};
\end{tikzpicture}}
\uline{\unsure{\deleted{C}c|\deleted{c}C}es deux trivialisations coïncident au dessus de la section unité} $1_P\times1_Q$ \uline{de} $P\times Q$, et définissent\inserted{\handwritten{donc}} une même section
\[
e_{E\deleted{\unsure{/S|/P|/Q}}}\handwritten{\in\Gamma}(E\handwritten{)}\deleted{\unsure{/S)|/P)|/Q)}}\ ,
\tag*{(2.2.3)}
\]
qu'on appellera par abus de langage la \uline{section unité} (ou \uline{section neutre}) de $E$. Pour prouver l'égalité annoncée\deleted{\unsure{, et une conséquence|et une conséquence}} des deux sections $e,e'$ de $E$ déduites respectivement de $e_{E/P}$ et $e_{E/Q}$, on note
\end{minipage}};
% Source-left marginal oval, alongside the opening partial-law line.
\path (draft.north west) -- (draft.south west) coordinate[pos=.084] (psiY);
\node[inner sep=0pt] at ([xshift=-1.1cm]psiY) {$\handwritten{\Psi}$};
\begin{scope}[shift={([xshift=-1.1cm]psiY)},rotate=-20]
\draw[gray,line width=.3pt] (0,0) ellipse[x radius=.29cm,y radius=.53cm];
\end{scope}
% Page-spanning cancellation geometry. Fractions keep the observed relative
% endpoints when the diplomatic prose reflows; no rectangular frame is added.
\path (draft.north west) -- (draft.south west) coordinate[pos=.227] (y227);
\path (draft.north west) -- (draft.south west) coordinate[pos=.891] (y891);
\path (draft.north west) -- (draft.south west) coordinate[pos=.008] (y008);
\path (draft.north west) -- (draft.south west) coordinate[pos=.685] (y685);
\path (draft.north west) -- (draft.south west) coordinate[pos=.083] (y083);
\path (draft.north west) -- (draft.south west) coordinate[pos=.815] (y815);
\path (draft.north west) -- (draft.south west) coordinate[pos=.426] (y426);
\path (draft.north west) -- (draft.south west) coordinate[pos=.658] (y658);
\path (draft.north west) -- (draft.south west) coordinate[pos=.975] (y975);
\draw[gray,line width=.32pt] ([xshift=-.09cm]y227) .. controls ([xshift=.3cm,yshift=-.7cm]y227) and ([xshift=1.8cm,yshift=.4cm]draft.south west) .. ([xshift=1.88cm,yshift=-.35cm]draft.south west);
\draw[gray,line width=.38pt] ([xshift=2.46cm,yshift=.27cm]draft.north west) .. controls ([xshift=2.8cm,yshift=-.5cm]draft.north west) and ([xshift=7.8cm]y685) .. ([xshift=9.17cm]y891);
\draw[gray,line width=.3pt] ([xshift=6.52cm]y008) -- ([xshift=11.52cm]y685);
\draw[line width=.43pt] ([xshift=11.29cm]y083) -- ([xshift=-.89cm]y815);
\draw[gray,line width=.38pt] ([xshift=12.23cm]y426) -- ([xshift=2.71cm,yshift=-.3cm]draft.south west);
\draw[gray,line width=.38pt] ([xshift=12.29cm]y658) .. controls ([xshift=12.7cm,yshift=-.3cm]y658) and ([xshift=10.8cm,yshift=.6cm]y975) .. ([xshift=9.03cm]y975);
\end{tikzpicture}}
\normalmarginpar
% archive sheet number: 26
